% !TEX root = ../../proposal.tex

\begin{abstract}
Compositional generative modeling aims to create complex unseen combinations by composing multiple pretrained models. 
Instead of relying on a single large-scale system, individual expert models can be combined at inference time to enable out-of-distribution generation. 
A single model must learn these combinations directly from training data, which becomes difficult when many of them are rare or missing.
In contrast, compositional approaches factorize this problem into reusable components, enabling data-efficient generalization to novel combinations without retraining.
This capability is critical in low-resource settings where specific data combinations are rare. 
In clinical imaging, for example, hospitals often have many scans for single diseases but very few examples of rare co-morbidities. 
A compositional system might bridge this gap by synthesizing plausible rare cases from these separate components to support data augmentation.

However, current compositional methods often fail. 
These systems struggle with tasks like attribute binding, negation, and combining closely related concepts. 
A primary reason for this failure is the mathematical operator used to combine conditions, which we term as ``logical composition''. 
This mathematical approach does not always match the semantically meaningful output produced by a single prompt for a complex scene. 
This creates a disconnect between the mathematical combination and the intended scene. 
The result limits image quality and makes synthetic data unreliable for sensitive clinical use.

The proposed study investigates this disconnect and aims to develop a method to reduce it. 
The research follows three steps. First, it characterizes the failure of logical composition through a structured categorization of concept pairs and trajectory analysis. 
Second, it tests whether learning more factorized representations, using a partitioned latent space, can align the mathematical combinations with the target scene. 
Finally, it demonstrates the utility of this framework in medical imaging. The proposed study evaluates whether rare, co-morbid chest X-rays can be synthesized from individual common and separate disease models. 
This work aims to make out-of-distribution synthesis more reliable and useful for practical downstream tasks.
\end{abstract}
